% Propietario conservado: /Users/ruben/Documents/New project/output/REVISION_ARGUMENTAL_APERTURAS_CIERRES_20260915/VII/delivery/MANUSCRIPT_VII_ES_EN_COMPLETE/source_es/manuscrito/31_corriente_y_respuesta_cartan.tex
% Artículo VII. Fragmento de cuerpo: corriente y respuesta Cartan--Holst.
% RESULTADO_RECUPERADO: acción, corriente variacional, inversa de Holst,
% torsión algebraica, respuesta cuadrática y normalización métrica.
% Fuentes leídas íntegramente:
% [II10] output/ARTICULO_II_REV10_EDICION_INTEGRADA_20260910/manuscrito/
%        sections/10_gravedad.tex, líneas 109--155.
% [G11c] integral activo 2249, fuente/manuscrito/sections/md/
%        11c_gravedad_holonomia_rev6.tex, líneas 214--334.
% [B028] integral activo, fuente/manuscrito/integracion_83/parte_iv/body/
%        B028_ch68_realizacion_cartan_holst_barbero_7df2afd9f675_11_gravedad_cartan_holst.tex,
%        líneas 76--218; copia idéntica conservada en II10/fuentes/propietarios_II/G03/.
%        SHA-256 e1e1f3325399b3ef09c331ed8268e3883f7bcc561cd7a560cf80d1be7eb1e91c.
% CERTIFICADO_NUEVO de resultados recuperados: inversión explícita del mapa
% torsión -> D(e wedge e), prueba de unicidad de contorsión y eliminación
% homogénea del término cuadrático. Las dos inversas se comprobaron sobre
% las 24 bases correspondientes mediante aritmética racional exacta.
% Precisión de dominio: la respuesta lineal en corriente y su eliminación
% cuadrática se enuncian para materia afín en la contorsión. Se conserva
% la dependencia del coeficiente respecto de la acción material concreta.
%
% DEPENDENCIAS DE INTEGRACIÓN (no constituyen remisiones bibliográficas):
% 1. vii:sec:realizacion-cartan: cotetrada, conexión y realización del
%    transporte enriquecido; incluir su prueba en el cuerpo precedente.
% 2. vii:sec:acoplamientos: genealogía de gamma_HMT, G_int, c_int,
%    homogeneidad dimensional y elección de la carta, con pruebas.
% 3. Para evaluar una especie material: acción S_m concreta, representación
%    de sus campos y cálculo de su corriente sigma. Este fragmento recibe
%    esa acción declarada y demuestra su respuesta variacional.
% 4. La identificación de una corriente de pantalla con sigma, su
%    contracción física y su normalización material tienen residencia propia.
%    Aquí no se introduce ni se selecciona una amplitud s_0.
% 5. El balance métrico utiliza la covariancia de la acción material y
%    sus ecuaciones de movimiento; no supone conservación separada del
%    tensor de espín y del tensor material de Levi--Civita.

\section{Corriente de espín y respuesta algebraica de Cartan--Holst}
\label{vii:vii:sec:corriente-respuesta-cartan}

La realización del transporte enriquecido de la
sección~\ref{vii:vii:sec:realizacion-cartan} proporciona la cotetrada y la
conexión; las secciones dimensionales y el parámetro de incidencia de la
secciones~\ref{vii:v:ant:sec:barbero} y~\ref{vii:v:ant:sec:gravity} fijan sus acoplamientos. Su composición
determina un problema variacional de primer orden. La corriente material
ocupa en él una posición precisa: es la derivada de la acción respecto de
la conexión y determina la respuesta torsional con la misma normalización.

\subsection{Dominio geométrico y normalización de la corriente}
\label{vii:vii:subsec:dominio-corriente-cartan}

Trabajamos en una variedad lisa orientada \(M\) de dimensión cuatro,
con fibrado lorentziano interno \(E\), métrica
\(\eta=\operatorname{diag}(-1,1,1,1)\) y cotetrada no degenerada
\(e\in\Omega^1(M;E)\). En un marco local escribimos
\(e^I\), con \(I=0,1,2,3\). La conexión métrica
\(\omega^{IJ}=-\omega^{JI}\) actúa mediante \(D_\omega\).
Cuando los campos \(\Psi\) sean espinoriales, se conserva una estructura
de espín y la representación material correspondiente. Definimos
\begin{equation}
 \Sigma^{IJ}=e^I\wedge e^J,\qquad
 F^{IJ}=\mathrm d\omega^{IJ}+\omega^I{}_{K}\wedge\omega^{KJ},
 \qquad T^I=D_\omega e^I.
 \label{vii:vii:eq:campos-cartan}
\end{equation}
El operador \(\star\) actúa sobre los índices bivectoriales internos,
\begin{equation}
 (\star X)^{IJ}=\tfrac12\epsilon^{IJ}{}_{KL}X^{KL},
 \qquad \star^2=-I,\qquad
 \mathsf P_\gamma=\star+\gamma^{-1}I.
 \label{vii:vii:eq:operador-holst}
\end{equation}
Se distingue así esta dualidad de la dualidad exterior de formas sobre
\(M\). Los índices internos se elevan y contraen con \(\eta\).

Fijamos una carta con \(G_{\rm int}>0\), \(c_{\rm int}>0\),
\(\gamma=\gamma_{\rm HMT}\in\mathbb R\setminus\{0\}\) y
\(\Lambda_{\rm int}\) constantes sobre ella. Abreviamos
\(\kappa=8\pi G_{\rm int}/c_{\rm int}^4\) y \(c=c_{\rm int}\).
Estas secciones permanecen fijas durante la variación. La acción es
\begin{equation}
 \begin{split}
 S[e,\omega,\Psi]={}&
 \frac1{2\kappa c}\int_M\Sigma_{IJ}\wedge
                         (\mathsf P_\gamma F)^{IJ}
 -\frac{\Lambda_{\rm int}}{\kappa c}
                         \int_M\operatorname{vol}_e
 +S_{\rm m}[e,\omega,\Psi].
 \end{split}
 \label{vii:vii:eq:accion-cartan-holst}
\end{equation}
La convención de volumen y orientación es la que da
\(\Sigma_{IJ}\wedge(\star F)^{IJ}=R\operatorname{vol}_e\)
en el sector de Levi--Civita. El tipo dimensional previo hace de cada
sumando una sección de acción.

Para la variación de la conexión se utilizan variaciones de soporte
compacto en el interior, o se fija
\(\delta\omega|_{\partial M}=0\). Puede emplearse también una
completación de borde cuyo problema variacional anule el mismo término
fronterizo. A cotetrada y campos materiales fijos, la corriente de espín
es la forma de grado tres, con valores en bivectores internos duales,
determinada por
\begin{equation}
 \delta_\omega S_{\rm m}
 =-\tfrac12\int_M\delta\omega^{IJ}\wedge\sigma_{IJ}.
 \label{vii:vii:eq:corriente-variacional}
\end{equation}
El emparejamiento con una variación arbitraria de soporte compacto fija
\(\sigma\) de manera única. Esta definición conserva la unidad de
acción de la corriente y el factor de normalización que interviene en
su respuesta geométrica.

\subsection{Variación de la conexión e inversión de Holst}
\label{vii:vii:subsec:variacion-holst}

\begin{theorem}[Ecuación variacional de Cartan--Holst]
\label{vii:vii:thm:ecuacion-cartan-holst}
En el dominio anterior, la estacionariedad respecto de la conexión
equivale a
\begin{equation}
 D_\omega(\mathsf P_\gamma\Sigma)^{IJ}
       =\kappa c\,\sigma^{IJ}.
 \label{vii:vii:eq:ecuacion-cartan-holst}
\end{equation}
\end{theorem}
\begin{proof}
La variación de la curvatura es
\(\delta F=D_\omega\delta\omega\). La dualidad interna es paralela
y autoadjunta para la contracción bivectorial; como \(\gamma\) es
constante, \(D_\omega\mathsf P_\gamma=\mathsf P_\gamma D_\omega\).
Poniendo \(B_{IJ}=(\mathsf P_\gamma\Sigma)_{IJ}\), se obtiene
\[
 B_{IJ}\wedge D_\omega\delta\omega^{IJ}
 =\mathrm d(B_{IJ}\wedge\delta\omega^{IJ})
   -D_\omega B_{IJ}\wedge\delta\omega^{IJ}.
\]
La última forma cambia de signo al permutar sus factores de grados tres
y uno. Las condiciones de borde eliminan la integral del término
exacto. El término cosmológico es independiente de \(\omega\).
Por tanto,
\[
 \delta_\omega S=
 \frac1{2\kappa c}\int_M\delta\omega^{IJ}\wedge
       \bigl[D_\omega B_{IJ}-\kappa c\,\sigma_{IJ}\bigr].
\]
La arbitrariedad de la variación prueba la ecuación, incluido su signo
y su coeficiente.
\end{proof}

\begin{lemma}[Inversa real del operador de Holst]
\label{vii:vii:lem:inversa-holst}
Para \(\gamma\in\mathbb R\setminus\{0\}\),
\begin{equation}
 \mathsf P_\gamma^{-1}
 =\frac{\gamma^2}{1+\gamma^2}
                      (-\star+\gamma^{-1}I).
 \label{vii:vii:eq:inversa-holst}
\end{equation}
\end{lemma}
\begin{proof}
Los dos factores son polinomios en \(\star\) y conmutan. Su producto
antes de normalizar es
\[
 (\star+\gamma^{-1}I)(-\star+\gamma^{-1}I)
 =-\star^2+\gamma^{-2}I=(1+\gamma^{-2})I.
\]
El denominador \(1+\gamma^2\) es positivo en el dominio real fijado.
La multiplicación por \(\gamma^2/(1+\gamma^2)\) da la identidad.
\end{proof}

La corriente determina así la forma bivectorial de grado tres
\begin{equation}
 Y^{IJ}:=\kappa c\,
   \frac{\gamma^2}{1+\gamma^2}
   (-\star+\gamma^{-1}I)\sigma^{IJ},
 \qquad D_\omega\Sigma^{IJ}=Y^{IJ}.
 \label{vii:vii:eq:corriente-dualizada}
\end{equation}

\subsection{Reconstrucción puntual de la torsión y de la contorsión}
\label{vii:vii:subsec:reconstruccion-torsion}

La regla de Leibniz da
\begin{equation}
 D_\omega\Sigma^{IJ}
 =T^I\wedge e^J-e^I\wedge T^J
 =:\mathsf L_e(T)^{IJ}.
 \label{vii:vii:eq:mapa-torsion}
\end{equation}
En cada punto, \(\mathsf L_e\) actúa de
\(\Lambda^2T^*M\otimes E\) en
\(\Lambda^3T^*M\otimes\Lambda^2E\). Ambos espacios tienen dimensión
veinticuatro. Sea \(E_I\) el marco dual de la cotetrada, de modo que
\(\iota_{E_I}e^J=\delta_I^J\).

\begin{theorem}[Inversión algebraica de la ecuación torsional]
\label{vii:vii:thm:inversion-torsion}
Para una cotetrada no degenerada, \(\mathsf L_e\) es un isomorfismo.
Su inversa sobre la forma \(Y\) se escribe
\begin{equation}
 B^I=\sum_J\iota_{E_J}Y^{IJ},\qquad
 b=\tfrac14\sum_I\iota_{E_I}B^I,\qquad
 T^I=B^I-e^I\wedge b.
 \label{vii:vii:eq:torsion-explicita}
\end{equation}
La ecuación~\eqref{vii:vii:eq:ecuacion-cartan-holst} determina por tanto
la torsión puntualmente a partir de \(e\) y de la corriente que figura
en ella.
\end{theorem}
\begin{proof}
Sea primero \(Y=\mathsf L_e(T)\) y pongamos
\(t=\sum_J\iota_{E_J}T^J\). La contracción de
\eqref{vii:vii:eq:mapa-torsion}, utilizando que \(T^I\) tiene grado dos
y que hay cuatro elementos en la cotetrada, da
\[
 \sum_J\iota_{E_J}(T^I\wedge e^J)=2T^I,
 \qquad
 \sum_J\iota_{E_J}(e^I\wedge T^J)=T^I-e^I\wedge t.
\]
De aquí \(B^I=T^I+e^I\wedge t\). Una segunda contracción produce
\[
 \sum_I\iota_{E_I}B^I=t+3t=4t.
\]
Se recupera \(t=b\) y después \(T^I=B^I-e^I\wedge b\).
En particular, el núcleo de \(\mathsf L_e\) es nulo. La igualdad
de dimensiones prueba su sobreyectividad y hace válida la fórmula
para toda forma \(Y\) del codominio.
\end{proof}

La conexión se descompone de manera única como
\begin{equation}
 \omega^{IJ}=\mathring\omega^{IJ}(e)+K^{IJ},\qquad
 K^{IJ}=-K^{JI},\qquad T^I=K^I{}_{J}\wedge e^J,
 \label{vii:vii:eq:contorsion}
\end{equation}
donde \(\mathring\omega(e)\) es la conexión de Levi--Civita y
\(K\) la contorsión. Para hacer explícita la inversa, escribimos
\(K_{IJ}=K_{IJK}e^K\) y
\(T_I=\tfrac12T_{IJK}e^J\wedge e^K\). Entonces
\begin{equation}
 T_{IJK}=K_{IKJ}-K_{IJK},\qquad
 K_{IJK}=\tfrac12(T_{KIJ}-T_{IJK}-T_{JKI}).
 \label{vii:vii:eq:contorsion-componentes}
\end{equation}
La segunda igualdad se obtiene combinando cíclicamente la primera y
usando \(K_{IJK}=-K_{JIK}\); su sustitución recupera las componentes
de \(T\). La torsión y la contorsión son, por consiguiente, dos
presentaciones equivalentes de la respuesta algebraica con cotetrada fija.

\begin{corollary}[Sector sin corriente y respuesta lineal mínima]
\label{vii:vii:cor:desacoplamiento-torsional}
Si \(\sigma=0\), se obtiene \(T=0\) y
\(\omega=\mathring\omega(e)\). Para un acoplamiento material cuya
corriente \(\sigma[e,\Psi]\) sea independiente de \(K\),
las ecuaciones~\eqref{vii:vii:eq:corriente-dualizada},
\eqref{vii:vii:eq:torsion-explicita} y
\eqref{vii:vii:eq:contorsion-componentes} proporcionan una respuesta única
\(K_*(e,\Psi)\), lineal en \(\sigma\).
\end{corollary}
\begin{proof}
Las tres aplicaciones de reconstrucción son lineales con \(e\),
\(\gamma\) y \(\kappa c\) fijos. La corriente nula produce
\(Y=0\), después \(T=0\) y finalmente \(K=0\).
\end{proof}

En la acción mínima, la ecuación de conexión contiene esta respuesta
puntual. Su evolución se obtiene al evolucionar \(e\) y \(\Psi\):
\(K_*=K_*(e,\Psi)\) cambia con ellos. Si una acción material conserva
dependencia algebraica adicional en \(K\), la ecuación de Cartan
mantiene su forma variacional y debe resolverse con esa corriente;
la respuesta lineal anterior corresponde al dominio mínimo especificado.

\subsection{Eliminación de la contorsión y contribución cuadrática}
\label{vii:vii:subsec:respuesta-cuadratica}

Consideremos ahora el acoplamiento mínimo afín en \(K\), expresado por
\begin{equation}
 S_{\rm m}[e,\mathring\omega+K,\Psi]
 =S_{\rm m}[e,\mathring\omega,\Psi]
   -\tfrac12\int_M K^{IJ}\wedge\sigma_{IJ}[e,\Psi].
 \label{vii:vii:eq:materia-afin-contorsion}
\end{equation}
Esta condición identifica la acción material a la que se aplica la
eliminación cuadrática; la corriente sigue siendo la definida por
\eqref{vii:vii:eq:corriente-variacional}.

La expansión de la curvatura es
\begin{equation}
 F(\mathring\omega+K)
 =F(\mathring\omega)+D_{\mathring\omega}K+K\wedge K,
 \qquad (K\wedge K)^{IJ}=K^I{}_{L}\wedge K^{LJ}.
 \label{vii:vii:eq:curvatura-contorsion}
\end{equation}
Como \(D_{\mathring\omega}\Sigma=0\), el término lineal es la
derivada exterior de \(\Sigma_{IJ}\wedge(\mathsf P_\gamma K)^{IJ}\).
La contribución de borde correspondiente se mantiene explícita:
\begin{equation}
 S_\partial[e,K]=\frac1{2\kappa c}
     \int_{\partial M}\Sigma_{IJ}\wedge(\mathsf P_\gamma K)^{IJ}.
 \label{vii:vii:eq:borde-contorsion}
\end{equation}
En un dominio interior, o con la completación de borde compatible,
la densidad de acción que depende de la contorsión es
\begin{equation}
 \mathcal Q_{e,\gamma}(K)-\tfrac12K^{IJ}\wedge\sigma_{IJ},
 \qquad
 \mathcal Q_{e,\gamma}(K)
  =\frac1{2\kappa c}\Sigma_{IJ}\wedge
                         \mathsf P_\gamma(K\wedge K)^{IJ}.
 \label{vii:vii:eq:forma-cuadratica-contorsion}
\end{equation}
La densidad \(\mathcal Q_{e,\gamma}\) es una forma diferencial de
grado cuatro, homogénea de grado dos en \(K\).

\begin{theorem}[Respuesta efectiva cuadrática en la corriente]
\label{vii:vii:thm:respuesta-cuadratica}
Bajo \eqref{vii:vii:eq:materia-afin-contorsion}, la sustitución de la
solución algebraica \(K_*\) produce la densidad efectiva de espín
\begin{equation}
 \mathcal L_{\rm spin}^{\rm eff}
 =-\tfrac14 K_*^{IJ}\wedge\sigma_{IJ}
 =-\mathcal Q_{e,\gamma}(K_*).
 \label{vii:vii:eq:accion-efectiva-spin}
\end{equation}
Con cotetrada y acoplamientos fijos, es homogénea de grado dos en
la corriente. La acción efectiva conserva además
\(S_\partial[e,K_*]\), o su completación de borde declarada.
\end{theorem}
\begin{proof}
Por el corolario~\ref{vii:vii:cor:desacoplamiento-torsional},
\(K_*\) es lineal en \(\sigma\). La ecuación de conexión es
precisamente la estacionariedad de
\eqref{vii:vii:eq:forma-cuadratica-contorsion}. Como ésta es algebraica,
se aplica puntualmente la variación radial \(\delta K=K_*\);
también puede multiplicarse por una función de soporte compacto.
La identidad de Euler para una forma cuadrática da
\[
 2\mathcal Q_{e,\gamma}(K_*)
             -\tfrac12K_*^{IJ}\wedge\sigma_{IJ}=0.
\]
La sustitución en la densidad de acción prueba
\eqref{vii:vii:eq:accion-efectiva-spin}. Si
\(\sigma\mapsto\lambda\sigma\), entonces
\(K_*\mapsto\lambda K_*\) y la densidad efectiva se multiplica
por \(\lambda^2\). La expansión
\eqref{vii:vii:eq:curvatura-contorsion} identifica separadamente el
término de borde que acompaña a esa eliminación.
\end{proof}

La corriente concreta, su contracción lorentziana, la dualidad y las
convenciones de la acción material determinan el signo y el coeficiente
de esta contribución. La identidad cuadrática proporciona la regla de
composición que permite evaluarlos una vez especificados esos datos.

\subsection{Fuente métrica efectiva y balance covariante}
\label{vii:vii:subsec:fuente-metrica-efectiva}

La normalización de la fuente métrica se fija en la misma recta de
acción. Tras eliminar \(K\), definimos
\begin{align}
 \delta_g S_{\rm m}[e,\mathring\omega,\Psi]
 &=-\tfrac12\int_M\operatorname{vol}_e\,
                   \mathcal T^{S,\rm LC}_{\mu\nu}\,\delta g^{\mu\nu},
 \\
 \delta_g S_{\rm spin}^{\rm eff}
 &=-\tfrac12\int_M\operatorname{vol}_e\,
                   \mathcal U^{S,\rm spin}_{\mu\nu}\,\delta g^{\mu\nu}.
 \label{vii:vii:eq:fuentes-normalizadas-accion}
\end{align}
Aquí se consideran variaciones interiores, con los términos de borde
tratados conforme al problema variacional fijado. Para el acoplamiento
afín, la segunda fuente es la variación métrica de
\eqref{vii:vii:eq:accion-efectiva-spin}. La variación incluye la dependencia
de la corriente material respecto de \(e\) y de \(\Psi\).

\begin{proposition}[Ecuación métrica y conservación de la fuente total]
\label{vii:vii:prop:fuente-metrica-balance}
Con las convenciones anteriores, las ecuaciones métricas son
\begin{equation}
 G_{\mu\nu}[g]+\Lambda_{\rm int}g_{\mu\nu}
 =\kappa c\bigl(\mathcal T^{S,\rm LC}_{\mu\nu}
                    +\mathcal U^{S,\rm spin}_{\mu\nu}\bigr)
 =\kappa\bigl(T^{\rm phys,LC}_{\mu\nu}
                    +U^{\rm phys,spin}_{\mu\nu}\bigr),
 \label{vii:vii:eq:ecuacion-metrica-efectiva}
\end{equation}
donde \(T^{\rm phys,LC}=c\mathcal T^{S,\rm LC}\) y
\(U^{\rm phys,spin}=c\mathcal U^{S,\rm spin}\) cuando la
cotetrada toma valores en la recta de longitud. Toda solución satisface
\begin{equation}
 \mathring\nabla^\mu
   (T^{\rm phys,LC}_{\mu\nu}+U^{\rm phys,spin}_{\mu\nu})=0.
 \label{vii:vii:eq:balance-metrico-total}
\end{equation}
\end{proposition}
\begin{proof}
La eliminación de \(K\) conserva las ecuaciones de las variables
restantes: en la regla de la cadena, el término que multiplica
\(\delta K_*\) se anula por la ecuación de conexión. En el sector
de Levi--Civita, la primera identidad de Bianchi anula el término
\(\Sigma_{IJ}\wedge F^{IJ}\), y la parte gravitatoria se escribe
\((2\kappa c)^{-1}\int_M(R-2\Lambda_{\rm int})\operatorname{vol}_e\).
Su variación interior es
\[
 \frac1{2\kappa c}\int_M\operatorname{vol}_e\,
        (G_{\mu\nu}+\Lambda_{\rm int}g_{\mu\nu})\delta g^{\mu\nu}.
\]
La suma con \eqref{vii:vii:eq:fuentes-normalizadas-accion} da
\eqref{vii:vii:eq:ecuacion-metrica-efectiva}. La identidad de Bianchi
contraída para \(\mathring\nabla\), la compatibilidad métrica y
la constancia de \(\kappa\) y \(\Lambda_{\rm int}\) en la carta
producen \eqref{vii:vii:eq:balance-metrico-total}.
\end{proof}

Antes de eliminar la conexión, las identidades geométricas conservan
la forma
\begin{equation}
 D_\omega T^I=F^I{}_{J}\wedge e^J,\qquad D_\omega F^{IJ}=0.
 \label{vii:vii:eq:bianchi-cartan}
\end{equation}
La primera es \(D_\omega^2e=F\wedge e\); la segunda se obtiene
expandiendo \(D_\omega(\mathrm d\omega+\omega\wedge\omega)\)
y cancelando los términos. En una acción material covariante, el
balance métrico total coincide con la identidad de Noether sobre las
ecuaciones materiales. Geometría y materia conservan así una fuente
conjunta; el reparto entre su contribución de Levi--Civita y su
contribución de espín queda determinado por la acción especializada.
